\section{Architecture}
\label{sec:method:architecture}

\todo{Appendix candidate: the departure-by-departure comparison against GNN+'s own equation
numbers within~\ref{sec:method:architecture:equations}, the itemised configuration list and
the paper-vs-reference-implementation asides in the normalisation and positional-encoding
paragraphs. Kept here: every equation the exactness and summarization claims
of~\ref{sec:method:architecture:summarization} and~\ref{sec:method:architecture:exact} cite
directly, since those claims cannot be checked against a bulleted configuration.}
The encoder follows GNN+ \citep{luo2024gnnplus}; hyperparameters come from the literature
and from the tuning sweep of~\ref{sec:method:experiment:hp}. One architecture is used for
unreduced, sparsified, partitioned and summarized graphs alike, with two exceptions:
partitioning changes the pooling path (\ref{sec:method:architecture:partitioning}), and the
exact-compression track changes the input schema (\ref{sec:method:architecture:exact}).

The encoder is an edge-aware graph convolutional network over AIGs. A node carries a one-hot
type indicator over $[\text{constant}, \text{PI}, \text{AND}, \text{PO}]$, so an AND gate is
$[0, 0, 1, 0]$, and an edge carries a one-hot indicator of whether it is inverted. The one
departure from that schema is the summarized track, where a node is a super-node and its
vector counts the members it absorbed: a super-node holding 2 primary inputs, 15 AND gates
and 1 primary output is $[0, 2, 15, 1]$, and a super-edge counts the inverted and
non-inverted wires merged into it. A one-hot vector is the count vector of a single-member
super-node, so both graph types occupy the same four-dimensional input space and are read by
the same encoder weights (\ref{sec:method:summary:schema}). The remaining settings are given
in~\ref{tab:method:architecture} and derived
in~\ref{sec:method:architecture:equations}.

\begin{table}[htbp]
\centering
\caption[Encoder configuration]{Encoder configuration, used unchanged for unreduced,
sparsified, partitioned and summarized graphs. The two exceptions are the pooling path
under partitioning and the input schema of the exact-compression track, described in
Sections~\ref{sec:method:architecture:partitioning}
and~\ref{sec:method:architecture:exact}.}
\label{tab:method:architecture}
\small
\begin{tabular}{ll}
\toprule
\textbf{Component} & \textbf{Setting} \\
\midrule
Node features & 4-dimensional type indicator (member counts when summarized) \\
Edge features & 2-dimensional inversion indicator \\
Positional encoding & Log logic level, projected to 32 dimensions \\
Input width & 160 (128-dimensional node projection concatenated with the encoding), \\
 & normalised over each graph before the first layer \\
Depth & 4 message-passing layers \\
Hidden width & 128, constant across layers \\
Activation & LeakyReLU in the encoder blocks, ReLU in the regression head \\
Feed-forward block & $128 \rightarrow 256 \rightarrow 128$ \\
Normalisation & Layer normalisation in graph mode, bracketing the feed-forward block \\
Residuals & Skip connection added after message passing in every layer \\
Dropout & $0.15$ in the encoder blocks, $0.3$ in the regression head \\
Degree normalisation & Disabled \\
Readout & Jumping knowledge summed over the 4 layers, then mean pooling \\
Regression head & $128 \rightarrow 64 \rightarrow 1$ with a terminal sigmoid \\
\bottomrule
\end{tabular}
\end{table}

\subsection{Encoder Formulation}
\label{sec:method:architecture:equations}
Nothing in this subsection is a contribution of this thesis: the formulation is GNN+'s
\citep{luo2024gnnplus}, written out in full because the exactness and summarization claims
of~\ref{sec:method:architecture:summarization} and~\ref{sec:method:architecture:exact} are
claims about \emph{this} forward pass and cannot be checked against a bulleted
configuration. The classic GCN layer it modifies \citep{kipf2017gcn} is
\begin{equation}
    h_v^{l} = \sigma\!\left( \sum_{u \in \mathcal{N}(v) \cup \{v\}}
    \frac{1}{\sqrt{\hat{d}_u \hat{d}_v}}\, h_u^{l-1} W^{l} \right).
    \label{eq:method:gcn}
\end{equation}
Throughout, $h_v^{l}$ is the representation of node $v$ after layer $l$, $e_{uv}$ the
attribute of edge $(u,v)$, $\sigma$ the activation, $\|$ concatenation, and biases are
suppressed for readability.

\paragraph{Edge-conditioned messages.}
Each edge attribute is projected by its own weight $W_e^{l}$ and fused with the neighbour
representation, so the message a node receives depends on the wire it arrived on:
\begin{equation}
    m_v^{l} = \sum_{u \in \mathcal{N}(v)}
    c_{uv} \, \sigma\!\left( h_u^{l-1} W^{l} + e_{uv} W_e^{l} \right).
    \label{eq:method:message}
\end{equation}
The per-edge coefficient $c_{uv}$ is the one place~\eqref{eq:method:message} is
configurable rather than fixed. With symmetric degree normalisation enabled it is the
GCN coefficient $1 / \sqrt{\hat{d}_u \hat{d}_v}$ of~\eqref{eq:method:gcn}, computed once
per edge when graphs are cached rather than at every forward pass; disabled, as it is
here, it is $1$ and no edge-weight tensor is carried at all. The same edge-weight
channel is what the exact-compression track of~\ref{sec:method:architecture:exact}
reuses: setting $c_{uv}$ to the edge multiplicity $w_{uv}$ scales the message
\emph{outside} $\sigma$, which is what makes one merged message equal the sum of the
messages it replaces.

\paragraph{Normalisation, dropout, residual and feed-forward block.}
GNN+ adds four of its six techniques inside the layer block:
\begin{align}
    \tilde{h}_v^{l} &= h_v^{l-1} + \operatorname{Drop}\!\left( \sigma(m_v^{l}) \right),
    \label{eq:method:residual} \\
    \bar{h}_v^{l} &= \operatorname{Norm}\!\left( \tilde{h}_v^{l} \right),
    \label{eq:method:norm1} \\
    h_v^{l} &= \operatorname{Norm}\!\left( \bar{h}_v^{l}
        + \operatorname{FFN}(\bar{h}_v^{l}) \right),
    \label{eq:method:layer} \\
    \operatorname{FFN}(h) &= \operatorname{Drop}\!\left(
        \operatorname{Drop}\!\left( \sigma(h W_1^{l}) \right) W_2^{l} \right),
    \label{eq:method:ffn}
\end{align}
with $W_1^{l} \colon 128 \rightarrow 256$ and $W_2^{l} \colon 256 \rightarrow 128$. GNN+
additionally normalises the aggregated message before the activation and the residual; that
step is dropped here, leaving only the pre- and post-FFN bracketing of GraphGPS
\citep{rampasek2022graphgps}. The two dropouts of~\eqref{eq:method:ffn} follow the GNN+
reference implementation.

\paragraph{Positional encoding.}
A per-node structural scalar, here the log-compressed logic level $\ell_v$, is
projected and concatenated to the node features rather than summed into them:
\begin{equation}
    h_v^{0} = \operatorname{GraphNorm}\!\left(
        x_v W_{\text{in}} \;\big\|\; \sigma(\ell_v W_{\text{PE}}) \right),
    \label{eq:method:input}
\end{equation}
with $x_v \in \mathbb{R}^{4}$, $W_{\text{in}} \colon 4 \rightarrow 128$ and
$W_{\text{PE}} \colon 1 \rightarrow 32$, so $h_v^{0} \in \mathbb{R}^{160}$. GNN+ instead
projects the joined vector back down to the model width. Keeping all $160$ dimensions
here means the first layer maps $160 \rightarrow 128$ and is the one layer whose residual
in~\eqref{eq:method:residual} is skipped, its input and output widths being unequal.

\paragraph{Configuration specific to this work.}
The encoder is GNN+; the following are this project's configuration of it, and none of
them change the technique:
\begin{itemize}
    \item \textbf{Degree normalisation off.} $c_{uv} = 1$ in~\eqref{eq:method:message},
    rather than the $1 / \sqrt{\hat{d}_u \hat{d}_v}$ of~\eqref{eq:method:gcn}. Fan-in in an
    AIG is fixed by node type, so that coefficient is constant within each type pair and
    adds nothing the type indicator does not already carry. GNN+'s own edge-aware layer
    disables it too, and its ablations do not cover the choice either way. The normalised
    variant remains reachable through the model configuration.
    \item \textbf{No self-loop.} Aggregation in~\eqref{eq:method:message} runs over
    $\mathcal{N}(v)$ without the $\cup \{v\}$ of~\eqref{eq:method:gcn}; a node's own
    state re-enters through the residual of~\eqref{eq:method:residual} only.
    \item \textbf{\texttt{LeakyReLU} for $\sigma$} in the encoder, and
    layer normalisation in graph mode for $\operatorname{Norm}$, taking statistics over all
    nodes of one graph, rather than the batch normalisation used throughout the GNN+
    ablations.
    \item \textbf{Level-based positional encoding} in~\eqref{eq:method:input}, a single
    learned projection of log logic depth, where GNN+ uses a Random Walk Structural
    Encoding (RWSE). Level comes from a single topological pass, where RWSE costs one
    sparse matrix product per walk step on graphs of this size and adds the walk length to
    a tuning budget already committed to the reduction parameters.
    \item \textbf{Readout and head.} GNN+ reads out from the last layer alone; here the
    layer outputs are summed first (jumping knowledge over all $L = 4$ layers), then
    mean-pooled and passed through a bounded head:
\end{itemize}
\begin{align}
    h_G &= \frac{1}{|\mathcal{V}|} \sum_{v \in \mathcal{V}} \sum_{l=1}^{L} h_v^{l},
    \label{eq:method:pooling} \\
    \hat{y} &= \operatorname{Sigmoid}\!\left(
        \operatorname{Drop}_{0.3}\!\left( \operatorname{ReLU}(h_G W_3) \right) W_4 \right),
    \label{eq:method:readout}
\end{align}
with $W_3 \colon 128 \rightarrow 64$ and $W_4 \colon 64 \rightarrow 1$. The terminal
sigmoid matches the $[0, 1]$ range of the optimizability target, and is the reason the
head uses \texttt{ReLU} and a heavier dropout ($p = 0.3$) than the $p = 0.15$ of the
encoder blocks.

\subsection{Partitioned Inputs}
\label{sec:method:architecture:partitioning}
Partitioning is the one reduction family that requires an architectural change. Graphs are
partitioned offline and the edges spanning partitions are removed, which leaves the node
set intact but disconnects the graph into $k$ components. Pooling directly to the graph
level would then average over components the encoder never allowed to communicate, so the
readout is made hierarchical:

\begin{itemize}
  \item GNN encoder (applied on intra-partition edges only)
  \item Mean pooling from node representations to partition representations
  \item Mean pooling from partition representations to a single graph representation
  \item Prediction head applied to the graph representation to produce the output $\hat{y}$
\end{itemize}

Feature normalization is also made partition-aware: the input \texttt{GraphNorm} normalizes
within each partition rather than within each graph, so statistics do not leak across a
boundary the message passing itself cannot cross.

\subsection{Sparsified Inputs}
\label{sec:method:architecture:sparsification}
Sparsification alters only the data pipeline: a sparsified AIG preserves the feature
schema of the original, so the encoder of~\ref{sec:method:architecture} ingests it
unmodified.

Reductions are stored as node and edge masks applied at load time rather than materialised
into the corpus, so a single corpus serves every sparsification setting. Edge masks (random
edge dropout, spanning forest) leave the node count unchanged, whereas node masks (PageRank
pruning, AND-gate-only) reduce it. The node budget governing batch assembly
(\ref{sec:method:experiment:batching}) is therefore measured on the post-mask count.

\subsection{Summarized Inputs}
\label{sec:method:architecture:summarization}
Summarized graphs are also fed to the unmodified encoder, which is possible only because
the merge preserves the input schema as a superset: node and edge features become member
counts rather than one-hot indicators, and a super-node containing a single node reproduces
its original feature row exactly (\ref{sec:method:summary:schema}). Positional encodings are
pooled by member minimum, so a super-node's level is that of its earliest member. This
adaptation has not yet been tested.
% TODO: unrun, confirm before defending positinal encoding adaptation

A full graph is consequently a valid input to a summary-trained model without preprocessing,
being a graph of size-1 super-nodes. The cross-state inference of RQ5 therefore evaluates
one set of weights on both representations.

\subsection{Encoder Variant for Exact Colour Refinement}
\label{sec:method:architecture:exact}
Colour refinement at count-cap $\infty$ (\ref{sec:method:summary:wl}) is exact rather than
approximate, but the guarantee holds only under a modified encoder.

When $k$ nodes merge, parallel edges collapse into one super-edge whose message must equal
the sum of the $k$ it replaces. Multiplicity $w_{uv}$ carried as an edge attribute enters
$\sigma$ in~\eqref{eq:method:message}, which does not decompose over sums. Setting
$c_{uv} = w_{uv}$ instead scales the message outside $\sigma$ and gives equality. Edge
attributes are therefore dropped, and inversion moves onto the nodes: fan-in is fixed by
node type in an AIG (zero for constants and primary inputs, two for AND gates, one for
primary outputs), so a node's type and its count of inverted incoming edges determine which
fan-ins are inverted, up to order. Which specific fan-in is inverted is lost, which a
graph-level target does not need.

Coarsening to depth $d$ equal to the number of message-passing layers then yields outputs
identical to the uncoarsened graph \citep{bollen2023exact}: compression that costs no
accuracy by construction, and the positive control for RQ5
(\ref{sec:method:experiment:design}). The variant pays for this in schema compatibility. A
full graph must be converted before this encoder can be queried on it, whereas the other
methods support cross-state inference with no such step.

Exactness (\ref{sec:prelim:reduction:exact}) needs two properties from the encoder: a
permutation-invariant sum over the neighbour multiset (\eqref{eq:prelim:exact-aggregation}),
and count-based node features. The former is the multiplicity coefficient $c_{uv} = w_{uv}$
established above. The latter is the summarized schema
of~\ref{sec:method:architecture:summarization}: it keeps the initial colouring able to
separate what the merge already separated, and supplies the member count the weighted
readout needs (\eqref{eq:prelim:exact-readout}).